% Generated by scripts/assemble.py; corrections belong in source/corrections.json.
\examhead{144}{中南大学 2024 年数学分析试卷}{137}
% p138-b05; source PDF page 138
\par\Needspace{5\baselineskip}
\blockanchor{p138-b05}
\majorquestion{一}计算题. 每题 10 分, 共 40 分.\par
% p138-b06; source PDF page 138
\blockanchor{p138-b06}
\question{1}求极限 $\lim_{n\to \infty}\left(1 - \frac{3}{4n^2}\right)^{n^2 + 4n + 3}.$\par
% p138-b07; source PDF page 138
\blockanchor{p138-b07}
\question{2}已知 $f(x) = \left| \begin{array}{ccc} e^x & x^2 & 2 \\ 1 & 3x & 4x^2 \\ 0 & 5 & e^{2x} \end{array} \right|$，求 $f'(x)$.\par
% p138-b08; source PDF page 138
\blockanchor{p138-b08}
\question{3}求定积分 $\int_0^{\frac{\pi}{2}}\frac{\cos\theta}{\cos\theta + \sin\theta}\mathrm{d}\theta .$\par
% p138-b09; source PDF page 138
\blockanchor{p138-b09}
\question{4}求级数 $\sum_{n=1}^{\infty} \frac{1}{n!}$ 的和.\par
% p138-b10; source PDF page 138
\par\Needspace{5\baselineskip}
\blockanchor{p138-b10}
\majorquestion{二}(10 分) 设 $\{f_n(x)\}$ 在 $[a, b]$ 上满足: 对任意的 $x \in [a, b]$, $f_n'(x)$ 存在. 且\par
% p138-b11; source PDF page 138
\blockanchor{p138-b11}
(1) 对任意的 $x_0 \in [a, b]$, $\{f_n(x_0)\}$ 收敛.\par
% p138-b12; source PDF page 138
\blockanchor{p138-b12}
(2) $\{f_n'(x)\}$ 在 $[a, b]$ 上一致收敛.\par
% p138-b13; source PDF page 138
\blockanchor{p138-b13}
证明： $\{f_n(x)\}$ 在 $[a,b]$ 上一致收敛.\par
% p138-b14; source PDF page 138
\par\Needspace{5\baselineskip}
\blockanchor{p138-b14}
\majorquestion{三}(20 分) 证明: $\int_{0}^{1} \frac{\ln(1 - t)}{t} \, \mathrm{d}t = -\sum_{n=0}^{\infty} \frac{1}{n^{2}}$.\par
% p138-b15; source PDF page 138
\par\Needspace{5\baselineskip}
\par\Needspace{6\baselineskip}
\blockanchor{p138-b15}
\majorquestion{四}(15 分) 设 $f(x)$ 在 $[0, \pi]$ 上二阶连续可导, $f(\pi) = 2$, 且\par
\nopagebreak[4]
% p138-b16; source PDF page 138
\blockanchor{p138-b16}
\displayformula{\int_ {0} ^ {\pi} [ f (x) + f ^ {\prime \prime} (x) ] \sin x \mathrm{d} x = 5.}
% p138-b17; source PDF page 138
\blockanchor{p138-b17}
求 $f(0)$\par
% p138-b18; source PDF page 138
\par\Needspace{5\baselineskip}
\par\Needspace{6\baselineskip}
\blockanchor{p138-b18}
\majorquestion{五}(15 分) 计算以下曲线积分\par
\nopagebreak[4]
% p138-b19; source PDF page 138
\blockanchor{p138-b19}
\displayformula{\int_ {C} x ^ {2} y \mathrm{d} x + (y - 3) \mathrm{d} y.}
% p138-b20; source PDF page 138
\par\Needspace{6\baselineskip}
\blockanchor{p138-b20}
其中 $C$ 为长方形区域的边界, 沿顺时针方向, 长方形四个顶点分别为 $(0,0), (0,1), (2,0), (2,1)$.
\majorquestion{六}(15 分) 定义以下二元函数\par
\nopagebreak[4]
% p138-b21; source PDF page 138
\blockanchor{p138-b21}
\displayformula{f: \mathbb {R} ^ {2} \to \mathbb {R}, f (x, y) = \left\{ \begin{array}{l l} \frac {x ^ {3} + x ^ {2} y + x y ^ {2} + 2 y ^ {4}}{x ^ {2} + y ^ {2}}, & x ^ {2} + y ^ {2} \neq 0; \\ 0, & x ^ {2} + y ^ {2} = 0. \end{array} \right.}
% p138-b22; source PDF page 138
\blockanchor{p138-b22}
证明：\par
% p138-b23; source PDF page 138
\blockanchor{p138-b23}
(1) $f$ 在 $(0,0)$ 点连续.\par
% p139-b00; source PDF page 139
\blockanchor{p139-b00}
(2) $f$ 在 $(0,0)$ 点处存在偏导数.\par
% p139-b01; source PDF page 139
\blockanchor{p139-b01}
(3) $f$ 在 $(0,0)$ 点处是否可微？请详细说明.\par
% p139-b02; source PDF page 139
\par\Needspace{5\baselineskip}
\blockanchor{p139-b02}
\majorquestion{七}(15 分) 设 $f: (0,1) \to \mathbb{R}$ 一阶连续可导, 且满足对任意的 $x, y \in (0,1)$, 有 $|f(x) - f(y)| \leq |x - y|^2$, 证明: $f$ 恒为常数.\par
% p139-b03; source PDF page 139
\par\Needspace{5\baselineskip}
\par\Needspace{6\baselineskip}
\blockanchor{p139-b03}
\majorquestion{八}(20 分) 已知 $f(x)$ 是 $(0, +\infty)$ 上的单调连续函数, 且 $\lim_{x \to +\infty} f(x) = 0$. 证明:\par
\nopagebreak[4]
% p139-b04; source PDF page 139
\blockanchor{p139-b04}
\displayformula{\lim _ {n \to \infty} \int_ {0} ^ {+ \infty} f (x) \sin (n x)   \mathrm{d} x = 0. (\text {可能有误})}
